\section{Terceiro objetivo - qualidade de serviço}

Neste objetivo iremos considerar uma implementação de serviços diferenciados,
em que cada nó possui 4 categorias de acesso, com 4 janelas mínimas de
contenção diferentes: $CW_{min0} = 20$, $CW_{min1} =~16$, $CW_{min2} =~10$ e
$CW_{min3} =~8$. Dentro de cada nó, se 2 categorias
quiserem aceder ao meio ao mesmo tempo, é dada prioridade à categoria mais
prioritária.

\vspace{\baselineskip}

Foi necessário fazer uma alteração na simulação para poder suportar os
serviços diferenciados. Para além disso, realizaram-se 10 simulações,
cada uma com 20000 iterações (ao contrário das 10000 usadas nos casos
abordados anteriormente) para aumentar a fiabilidade dos resultados.
Os valores utilizados na simulação foram 10 nós e \textit{m} = 8.
Os tamanhos dos slots ultilizados foram $T_{idle} = 1~u.t.$,
$T_{succ} = 150~u.t.$ e $T_{coll} = 160~u.t.$

\vspace{\baselineskip}
\begin{table}[h]
    \centering
    \begin{tabular}{c|c|c|c}
        $CW_{min0}$ & $CW_{min1}$ & $CW_{min2}$ & $CW_{min3}$\\
        \midrule
        13,75\% & 16,63\% & 28,29\% & 35,37\%
    \end{tabular}
    \caption{Valores de débito para cada categoria de acesso}
    \label{tab:objetivo3}
\end{table}
\vspace{\baselineskip}

Os resultados do débito para cada categoria de acesso podem ser observados
na Tabela \ref{tab:objetivo3}. Podemos observar que as categorias com maior
prioridade beneficiaram de valores de débito mais altos, como esperado.
Mesmo assim, as categorias de menor prioridade não foram abafadas pelas de
maior prioridade, e conseguiram débitos razoáveis. É possível que, para
outros números de nós na rede e valores de \textit{m}, os resultados não
fossem tão positivos, por razões já discutidas anteriormente.